% =====================================================================
% Appendix: Formal Verification Plan.
%
% The plan that carries out ADR-VERIFICATION.FORMAL-01 and
% ADR-VERIFICATION.MODELS-01: the checkers and what each reaches, the
% properties each one checks, the phases of the work, the Lean
% libraries, when the checks run, and what stays unchecked. The decisions are in those records; a property listed here
% becomes normative with the test case that checks it.
% =====================================================================

This appendix is the plan that carries out
\reqref{ADR-VERIFICATION.FORMAL-01} and \reqref{ADR-VERIFICATION.MODELS-01},
which answer \reqref{HLR-VERIFICATION.FORMAL-01} and
\reqref{HLR-VERIFICATION.MODELS-01}. The decisions are in those records;
this appendix lists what each check is to establish, and the order the work
takes. A property listed here becomes normative when the test case that
checks it is written: each phase writes the test cases of its checks, or new
versions of the test cases its checks extend. What it says of the tools was
checked in November 2026. They change quickly, and each phase begins by
checking again what it relies on.

\section{The Checkers}
\label{app:verification-checkers}

The checkers are Kani,\footnote{\url{https://github.com/model-checking/kani};
its function contracts and its stubbing are unstable features, enabled with
\texttt{-Z function-contracts} and \texttt{-Z stubbing}.}
loom,\footnote{\url{https://github.com/tokio-rs/loom}}
shuttle,\footnote{\url{https://github.com/awslabs/shuttle}}
Veil,\footnote{\url{https://github.com/verse-lab/veil}} and Aeneas with its
front end Charon,\footnote{\url{https://github.com/AeneasVerif/aeneas} and
\url{https://github.com/AeneasVerif/charon}.} beside the replay of a
device-mapper \texttt{log-writes} device that
\reqref{TC-PERSISTENCE.CRASHSAFE-02} already specifies.
Table~\ref{tab:verification-checkers} gives what each establishes and its
limits, as of November 2026.

\begin{table}[tb]
\centering\small\renewcommand{\arraystretch}{1.3}%
\caption{The checkers}\label{tab:verification-checkers}
\begin{tabularx}{\linewidth}{|>{\raggedright\arraybackslash}p{0.15\linewidth}|>{\raggedright\arraybackslash}X|>{\raggedright\arraybackslash}p{0.17\linewidth}|>{\raggedright\arraybackslash}X|}
\hline
\textbf{Checker} & \textbf{Establishes} & \textbf{On the code itself} & \textbf{Limits} \\
\hline
Kani & That no input up to a bound makes the code panic, overflow, or fail an assertion & Yes, \texttt{unsafe} code included & One thread, no \texttt{await}; system calls stubbed; contracts and stubbing unstable \\
\hline
loom 0.7 & That a small concurrent test passes under every interleaving of its threads, up to a bound on preemptions & Yes, its \texttt{std::sync} types swapped & No Tokio; by its own documentation, misses some executions the C11 memory model allows \\
\hline
shuttle & That a concurrent test passes under many schedules, random or chosen by probabilistic concurrency testing & Yes, Tokio included through its wrappers & A test, not a proof \\
\hline
Veil 2.0 & That the invariants of a model written in Lean hold: by SMT solvers where they can, by Lean proof where they cannot; small instances also checked exhaustively & No: a model & Pins Lean 4.35.0 \\
\hline
Aeneas and Charon & Any theorem about the Rust code, as translated into Lean & A subset: no \texttt{async}, \texttt{unsafe}, trait objects or \texttt{Drop} behaviour & Charon in beta; the Lean library needs Lean 4.36.0 and Mathlib \\
\hline
\texttt{log-writes} replay & The disk state after each write of a real run, and after each block of a longer one, the flush and FUA points among them & Yes, on the real kernel & Root and a scratch device; one kernel and one file system per run; the writes in the order the device completed them \\
\hline
Exhaustive test & That a property holds in every case of a domain small enough to go through, one case at a time & Yes & Only where the cases are few: the orders of a few events, the indices of a small table \\
\hline
\end{tabularx}
\end{table}

\section{What Is Checked}
\label{app:verification-properties}

Each property below names the requirements it checks. The data formats, the
\texttt{unsafe} code and the concurrency are checked on the code itself; the
protocol is checked on a model, tied to the code by the tests' logs, and in
the end by proofs about the code.

\subsection{Data formats}

With Kani, for every input up to a bound; in phase 5, with proofs in Lean
for inputs of every size.
\begin{itemize}
  \item \emph{Key segment}: decoding an encoded key gives the key back, no
    two keys share an encoding, and decoding any bytes returns a key or an
    error, never a panic (\reqref{SDS-LAYOUT.KEYSEGMENT-01}).
  \item \emph{Record header}: parsing a built header gives back its fields,
    its sealed CRC matches, and parsing any 44 bytes returns a header or an
    error, never a panic or an overflow (\reqref{SDS-LAYOUT.OBJECTFILE-03}).
    The header has a fixed size, so the check covers every possible header,
    not a sample.
  \item \emph{Read-time validation}: a record is accepted only when every
    check of \reqref{SDS-LAYOUT.OBJECTFILE-03} passes, and malformed bytes
    never make the reader panic.
  \item \emph{Sharding}: an object's path depends only on its key's combined
    hash and tag, stays inside the cache directory, and ends in the 32
    hexadecimal characters of its name (\reqref{SDS-LAYOUT.SHARDING-01}). The
    hash is stubbed, since the property holds whatever its value.
  \item \emph{Shard-directory bitset}: every index it computes is in range
    (\reqref{FDS-LAYOUT.STARTUP-01}).
  \item \emph{Expiry never lengthens the metadata}: for every stored deadline
    and every instant, the internal-metadata segment an expiry writes is no
    longer than the one it replaces, through the MessagePack encoder of
    \texttt{pingora-cache} itself (\reqref{FDS-STORAGE.EXPIRE-01}). Should
    Kani prove unable to check the encoder, the existing tests, which cover every
    width the encoding gives a deadline, remain the check.
  \item \emph{Writer-lease names}: parsing a lease's name gives back the
    identity it was formatted from (\reqref{ADR-STORAGE.WRITERID-01}).
\end{itemize}

\subsection{Unsafe code}

The reads with Kani, their system calls stood in for by every result they
may return; the ring's table of operations in flight with an exhaustive
check of every order of its events, as Kani runs out of memory on it.
\begin{itemize}
  \item \emph{Reads that must not wait}: for every result \texttt{preadv2}
    or \texttt{pread} may return, the buffer's new length covers only bytes
    the call wrote (\reqref{FDS-STORAGE.PAGECACHE-01}).
  \item \emph{The ring}: a buffer submitted to the \texttt{io\_uring} ring,
    and the file it writes to, stay alive until the kernel's completion for
    it is collected, for any sequence of submissions, completions and drops
    (\reqref{FDS-STORAGE.IOURING-01}).
\end{itemize}
The other \texttt{unsafe} blocks call \texttt{libc} functions that take no
pointer --- \texttt{flock} on a descriptor the caller borrows,
\texttt{sync}, \texttt{geteuid} --- and leave nothing to check beyond the
call.

\subsection{Concurrency}

With loom where the code uses only \texttt{std::sync}, and with shuttle where
Tokio is involved and shuttle can run it (see phase~1 below: none so far).
\begin{itemize}
  \item \emph{Publication lock and read-only gate} (loom): one change per key
    at a time; when the switch to read-only returns, every change that
    entered the gate before it has ended, and none enters afterwards; no
    deadlock (\reqref{FDS-COMPAT.READONLY-02}).
  \item \emph{Inline slots} (loom): never more slots taken than the limit,
    nor than the runtime's workers less one, and every operation gives its
    slot back, on every path out of it
    (\reqref{FDS-STORAGE.INLINEWRITES-02}).
  \item \emph{In-flight registry} (loom): a purge that invalidates a key
    ends the key's in-flight stores, the readers waiting on them wake, and a
    store's entry leaves the registry when the store completes or aborts
    (\reqref{FDS-STORAGE.PARTIAL-02}, \reqref{FRE-COMPLIANCE.PURGE-01}).
\end{itemize}
That the store waiting on each write the ring takes learns that write's
result exactly once, whatever order the completions arrive in, needs no
checker of threads: the ring's table is used under its mutex, so its events
come one at a time, and the exhaustive check of the ring's table goes
through every order of them.

\subsection{The protocol across crashes}

Invariants of the protocol model, which Veil proves for any number of keys,
writers and calls --- or a Lean proof does, where its solvers cannot --- and
which the crash explorer of phase 3 checks on the real backend:
\begin{enumerate}
  \item \emph{Synchronised stores}: a published name never refers to data
    that is not on stable storage (\reqref{ADR-PERSISTENCE.DURABILITY-02},
    \reqref{FDS-STORAGE.PUBLISH-03}, \reqref{ADR-STORAGE.ATOMICPUBLISH-01}).
  \item \emph{No mixtures}: in every state a crash can leave, the restarted
    backend serves each object as one version it stored, or misses it ---
    after stores, in-place metadata updates and withdrawals alike
    (\reqref{HLR-PERSISTENCE.CRASHSAFE-02},
    \reqref{FDS-STORAGE.METAUPDATE-01}).
  \item \emph{Purges hold}: no crash undoes a purge the backend acknowledged
    (\reqref{FRE-COMPLIANCE.PURGEDURABLE-01}).
  \item \emph{Purges end in-flight stores}: no store in flight when a purge
    invalidates its key lands after the purge
    (\reqref{FRE-COMPLIANCE.PURGE-01}, \reqref{FDS-STORAGE.PURGE-01}).
  \item \emph{Unsynchronised stores} (\reqref{ADR-PERSISTENCE.UNSYNCED-01},
    \reqref{FDS-PERSISTENCE.CRASHDISCARD-01}): with the discard on, no record
    or temporary written before a host restart is reachable after it; with
    the discard off, on a file system that records a file's length only once
    its data is written, a record whose body did not reach the device reads
    short, and is a miss.
  \item \emph{The restart scan removes only leftovers}: never the temporary
    of a writer still running, process identifiers reused and hot upgrades
    included (\reqref{FDS-STORAGE.RECLAIM-01},
    \reqref{ADR-STORAGE.WRITERID-01}).
  \item \emph{Recovery beside live traffic} admits no ghost entry, one for a
    record that is gone, and leaves the eviction manager's entries matching
    the records that are published (\reqref{FDS-EVICTION.STARTUP-01}).
  \item \emph{Hot upgrade}: once the switch to read-only has returned, the
    old process publishes, purges, updates, expires and reclaims nothing
    (\reqref{FDS-COMPAT.READONLY-02}, \reqref{FRE-COMPAT.HOTUPGRADE-01}); its
    flush still removes the markers of unsynchronised stores
    (\reqref{FDS-PERSISTENCE.CRASHDISCARD-01}).
\end{enumerate}

\subsection{The code against the model}

In phase 5, proofs in Lean that the synchronous code does what the protocol
model's steps say: publishing a store, purging, expiring, updating
metadata, and classifying what the restart scan finds. Each proof assumes
that the caller holds the key's publication lock, which the concurrency
checks and the model establish.

\section{How the Checks Are Kept Honest}
\label{app:verification-rules}

The rules of \reqref{ADR-VERIFICATION.FORMAL-01}, in practice:
\begin{itemize}
  \item Every run lists the axioms each theorem rests on (Lean's
    \texttt{\#print axioms}) and fails on one outside the reviewed list, and
    a proof that uses \texttt{sorry} fails the build. A proof that needs
    more than Lean's default allowance of work for one proof
    (\texttt{maxHeartbeats}) raises it for itself alone, where a reviewer
    sees it; that is no limit on how long the check runs.
  \item Each Kani stub, and each model a proof assumes of a dependency, is
    named in the same list.
  \item Kani and loom harnesses are named \texttt{tc\_<id>\_\ldots}, after
    the test case they implement, so that \texttt{cargo xtask trace-report}
    and the tracing report count them as they count tests. A Lean file that
    checks a test case carries the marker
    \texttt{-{}- IMP-TC.<ID>: satisfy TC-<ID>}.
  \item A change to checked code updates its checks in the same pull
    request. Lake, the Lean build tool, checks again only the files a change
    reaches and the files that depend on them.
\end{itemize}

\section{Phases}
\label{app:verification-phases}

Each phase is useful on its own. Phases~0 and~1 need no Lean code: they
prepare the ground and check the code itself. Each later phase builds on
what the earlier ones established.
The pilot of phase~4 measures what the proofs about the code cost before
phase~5 commits to them.

\subsection*{Phase 0: groundwork}
\begin{itemize}
  \item The two tools that link test cases to code,
    \texttt{reqman-report.py} and \texttt{cargo xtask trace-report}, read
    Rust files only; they learn to read the markers in Lean's
    \texttt{-{}-} comments.
  \item The header check (\texttt{cargo xtask headers}) covers Lean files,
    and accepts in each library's directory that library's own copyright
    line.
  \item The crate declares two configuration names to the compiler's check
    of them: \texttt{kani}, which Kani sets, and
    \texttt{durable\_cache\_loom}, which the loom checks build the crate's
    tests with. They do not use the customary \texttt{loom}: a
    \texttt{-{}-cfg} reaches every crate a build compiles, and a dependency,
    \texttt{event-listener}, does not build as one under \texttt{cfg(loom)}.
  \item A module \texttt{crate::sync} gives the code the \texttt{std::sync}
    types, or loom's in the tests built with
    \texttt{cfg(durable\_cache\_loom)}, one declaration with a
    \texttt{\#[cfg]} for each, as \texttt{crate::platform} is laid out. loom
    never poisons a lock, so there its locks gain the two methods of the
    standard library's that it leaves out, and no lock is ever poisoned.
  \item The Nix shell provides elan, the Lean toolchain manager. Kani is not
    in nixpkgs: it installs with \texttt{cargo install kani-verifier} and
    \texttt{cargo kani setup}, and in CI with its GitHub action. Charon and
    Aeneas join in phase~4.
  \item A \texttt{verification/} directory holds the record of checked
    properties, and the CI workflows hold the jobs of
    Section~\ref{app:verification-schedule}, none with a time limit of its
    own.
  \item The two designs \reqref{ADR-VERIFICATION.FORMAL-01} expects --- how
    the checks run, and where they live --- are written.
\end{itemize}

\subsection*{Phase 1: the code itself}
\begin{itemize}
  \item Kani harnesses for the data formats and the \texttt{unsafe} code,
    each with small bounds in the run on each change, and deeper bounds in
    the daily run.
  \item loom tests for the publication lock, the read-only gate and the
    inline slots; shuttle tests for the in-flight registry and the ring's
    completions.
  \item The \texttt{log-writes} harness \reqref{TC-PERSISTENCE.CRASHSAFE-02}
    specifies, which the repository does not have yet: the test case's test
    builds, by hand, the states such a replay produces. The harness runs
    daily in a privileged job, on ext4 as the test case says, and on XFS too.
    Should the kernel of the CI runners lack the device-mapper target, it
    needs a runner of the project's own, or a virtual machine.
  \item Exit: every check passes, with a test case for each property, and
    the record of checked properties shows how long each check's runs take.
\end{itemize}
Phase~1 is carried out in parts, the first of which brings the runner of
the checks (\reqref{FDS-VERIFICATION.RUNS-01}), the loom checks of the
publication locks, the read-only gate and the inline slots
(\reqref{TC-STORAGE.CHANGELOCKS-01}, \reqref{TC-STORAGE.INLINESLOTS-01}),
and the Kani checks of the key segment and the record header
(\reqref{TC-LAYOUT.FORMATS-01}). What it found:
\begin{itemize}
  \item loom bounds nothing unless told to, and three threads that each take
    a slot then run for hours. The checks allow at most two preemptions on
    each change and three daily, which keeps every one under 15~seconds.
  \item Kani cannot follow \texttt{std::str::from\_utf8} over bytes whose
    values or addresses it does not know: the loop that validates UTF-8 a
    word at a time has no bound it can find. The key segment's UTF-8 check
    is therefore one function of its own, and Kani's checks stand in for it
    with one that accepts the ASCII tags, which the record lists as an
    assumption. For the same reason the round trip is checked for every
    primary hash and variance but fixed user tags, and decoding for every
    byte string of up to 37~bytes.
  \item \texttt{crc32c} chooses its implementation by asking the processor,
    which Kani cannot follow either; the header's checks replace it with
    another function of the same bytes, as their properties hold for any.
\end{itemize}
The second part brings the \texttt{unsafe} code
(\reqref{TC-STORAGE.UNSAFE-01}) and the shards of names
(\reqref{TC-LAYOUT.NAMES-01}), and a third kind of check: the exhaustive
check, an ordinary test that goes through every case where the cases are
few enough and Kani cannot finish. What it found:
\begin{itemize}
  \item Kani cannot stand in for a function of another language, such as
    \texttt{libc::pread64}, nor for a method of a type of the standard
    library, such as \texttt{io::Error::last\_os\_error}. The read that waits
    for the device is therefore a function of the crate's own,
    \texttt{pread\_blocking}, beside \texttt{pread\_nowait}, and the stand-in
    for both sets \texttt{errno}, as \texttt{libc} does.
  \item Kani walks every branch it cannot rule out, a formatted error
    message included, and runs out of memory there; the check of the reads
    stands in for the function that formats one.
  \item Kani runs out of memory on the ring's table, a vector of slots
    whose states it follows through every event, even for one operation. The
    table's three transitions are now methods of the table --- a completion,
    a poll, a drop --- which an exhaustive check drives through each of
    10,735 orders of two operations' events.
  \item Kani runs out of memory, too, on the shard-directory bitset's 1024
    words, which an exhaustive check goes through shard by shard, and on
    four properties that therefore stay with the tests, as the record of
    checked properties says: the read-time validation of a record, whose
    buffers have lengths read from the header; the shape of an object's
    name, built a character at a time; the names of writer leases,
    formatted through \texttt{std::fmt}; and the expiry's length, as the
    encoder of \texttt{pingora-cache} compresses the header segment with
    zstd, C code beyond Kani.
\end{itemize}
The third part brings the in-flight registry
(\reqref{TC-STORAGE.INFLIGHT-01}), with loom rather than shuttle: the
registry is a lock, atomics and a \texttt{OnceLock}, with no Tokio in it, so
\texttt{crate::sync} gives it loom's types as it does the publication locks.
Where Tokio is involved --- a partial hit's reader waiting on Tokio's
\texttt{Notify} for the next chunk --- the wait sits in the hit handler's
read loop, beside reads on the blocking pool, which neither loom nor
shuttle runs without the runtime; it stays with the streaming tests, as
the record of checked properties says, and phase~1 uses no shuttle check.
The fourth part brings the \texttt{log-writes} harness,
\texttt{cargo xtask crash-replay}. The kernel of GitHub's runners has the
device-mapper target, so the harness runs there, as root, daily and by
hand. It runs the scenario of \reqref{TC-PERSISTENCE.CRASHSAFE-02} on a file
system over the logged device, marking the log after each operation the
backend acknowledges; replays the log, with a replay of its own over the
format the kernel writes; and restarts the backend on each state the replay
reaches. What it found:
\begin{itemize}
  \item At the flush and FUA points the test case names, a store published
    without its synchronisation is never served wrong: it is a miss, which
    a host crash is allowed to cause. ext4 in its default
    \texttt{data=ordered} mode and XFS write a file's data before the
    metadata that gives the file its length, and ext4 in
    \texttt{data=writeback} mode, with delayed allocation, allocates a
    file's blocks and records its length only when it writes the data back.
    With the backend's \texttt{fdatasync} taken out, every one of those
    replays passed.
  \item The harness therefore checks every state, not only those points:
    after each write, and after each 4~KiB block of a longer one, as a power
    failure may leave a write torn. It also mounts ext4 with
    \texttt{data=writeback} and without delayed allocation, which records a
    file's length with the next commit of its journal, whether the data is
    written or not. There, with the \texttt{fdatasync} taken out, a write
    torn after a record's header serves the header over a body never
    written, and the check fails; with the directory's \texttt{fsync} after
    a purge taken out, a purged object is served on each file system. With
    the code as it is, the five runs pass: 4,082 states, about 800 each, in
    7~minutes.
  \item The states are those of the writes the device completed, in the
    order it completed them. A device with a volatile write cache may also
    lose some writes made before a flush and keep later ones; listing those
    states is the file-system library's task, in phase~3.
\end{itemize}
With this part, phase~1 is complete: every check passes, and the record of
checked properties gives the duration of each.

\subsection*{Phase 2: the protocol in Lean}
\begin{itemize}
  \item The storage-contract library, a first version: what
    \texttt{pingora-cache} expects of \texttt{Storage}, \texttt{HandleHit},
    \texttt{HandleMiss} and \texttt{EvictionManager}, with a reference
    implementation that runs. \texttt{MemCache}, the backend
    \reqref{TC-COMPLIANCE.RFC-02} compares against, is run through it:
    everything it does must be allowed.
  \item The protocol model, in Veil. Its state: names, records, temporaries,
    leases, the per-key locks, the read-only gate, the in-flight stores, and
    which data and directory entries are on stable storage. Its steps: the
    storage calls, a crash, a restart, the restart scan, a hot-upgrade
    handover. The invariants of Section~\ref{app:verification-properties}
    are proved, and small instances also checked exhaustively.
  \item Every behaviour of the model is one the storage contract allows.
  \item Trace checking: the test hooks, which compile to nothing outside
    tests, log the steps the backend takes; a Lean program replays each test
    run's log through the model, and fails on a step the model forbids.
  \item Exit: the invariants proved, and the logs of the crate's tests and
    of the compatibility tests replayed without a forbidden step.
\end{itemize}

\subsection*{Phase 3: the file system and its crashes}
\begin{itemize}
  \item The file-system library: the POSIX calls the backend makes; what a
    crash may leave of them under four profiles --- the POSIX minimum, ext4
    with \texttt{data=ordered}, ext4 with \texttt{data=writeback}, and XFS;
    a program that lists every state a crash may leave after a sequence of
    calls; and a harness that compares its behaviour without crashes with
    Linux.
  \item The crash explorer, on the method of
    ALICE:\footnote{T.~S.~Pillai, V.~Chidambaram, R.~Alagappan, S.~Al-Kiswany,
    A.~C.~Arpaci-Dusseau and R.~H.~Arpaci-Dusseau, ``All File Systems Are Not
    Created Equal: On the Complexity of Crafting Crash-Consistent
    Applications'', \emph{OSDI~2014}.} record the system calls of the
    scenario tests with \texttt{strace}, which the benchmark's census of
    system calls already uses; list the states the library allows a crash to
    leave;
    build each state in a scratch directory and restart the backend on it;
    and check the expected result of \reqref{TC-PERSISTENCE.CRASHSAFE-02}.
    Small scenarios run on each change, every scenario of the compatibility
    tests daily.
  \item The cross-check: each state a replay of the \texttt{log-writes}
    harness produces must be one the library's profile for that file system
    allows, which checks the library against the kernel.
\end{itemize}

\subsection*{Phase 4: the pilot of Aeneas}
\begin{itemize}
  \item Translate the key codec, then the sharding, then the record format,
    with the crate models.
  \item Known obstacles. The guard that edits a record header reseals the
    header's CRC when it is dropped, and Charon does not model what
    \texttt{Drop} does; an explicit call in its place would make the code
    worse, so that proof waits for Charon. The restart scan takes the
    eviction manager as a trait object, which Charon does not translate.
    Error paths and logging become calls the proofs treat as opaque.
  \item Measured: what translates unchanged, the size of the proofs, the
    time they take to check, and how much proof a typical change to the code
    breaks.
  \item Criteria: no code made worse for the translation, and little proof to
    repair after a typical change. The time the proofs take to check is
    measured and recorded, for the later decision on how long a check may
    run.
    Should the pilot miss one, a revision of
    \reqref{ADR-VERIFICATION.FORMAL-01} decides how phase~5 proceeds: with
    Aeneas once the Charon issues the pilot hit are closed, or with Verus.
\end{itemize}

\subsection*{Phase 5: proofs of the synchronous code}
\begin{itemize}
  \item The library of Rust Unix I/O, built on the file-system library.
    Every outcome a call may have --- a short read, \texttt{EINTR},
    \texttt{EAGAIN}, an error --- is an input the theorems quantify over.
  \item Module by module, as each settles: publishing a store, purging,
    expiring, updating metadata, classifying what the restart scan finds;
    and the data-format properties of phase~1 for inputs of every size. A
    module joins the record of checked properties when its proofs are
    complete.
\end{itemize}

\section{The Lean Libraries}
\label{app:verification-libraries}

Table~\ref{tab:verification-libraries} lists the four libraries of
\reqref{ADR-VERIFICATION.MODELS-01} by what they model; each is named when it
is created.

\begin{table}[tb]
\centering\small\renewcommand{\arraystretch}{1.3}%
\caption{The Lean libraries}\label{tab:verification-libraries}
\begin{tabularx}{\linewidth}{|>{\raggedright\arraybackslash}p{0.15\linewidth}|>{\raggedright\arraybackslash}X|>{\raggedright\arraybackslash}p{0.17\linewidth}|>{\raggedright\arraybackslash}p{0.2\linewidth}|}
\hline
\textbf{Library} & \textbf{Models} & \textbf{Depends on} & \textbf{Checked against} \\
\hline
File system & The POSIX calls the backend makes; what a crash may leave of them, per profile; the list of those states & Lean core & Linux; the \texttt{log-writes} replays \\
\hline
Rust Unix I/O & \texttt{std::fs}, \texttt{std::io} and \texttt{std::os::unix}; \texttt{libc}; \texttt{rustix} & The Aeneas library; the file-system library & The real calls \\
\hline
Crates & \texttt{bytes}; \texttt{hex}; \texttt{zerocopy}, its little-endian integers and \texttt{repr(C)} layout; \texttt{crc32c}, with the law for extending a running CRC that a record's body CRC relies on; \texttt{blake2}, as an arbitrary function; \texttt{thiserror}; \texttt{log}; \texttt{pingora-error} & The Aeneas library & Property tests against the crates \\
\hline
Storage contract & The types of \texttt{pingora-cache}, the contract of its storage traits, the laws of \texttt{CacheMeta} & Lean core & \texttt{MemCache}, through the reference implementation \\
\hline
\end{tabularx}
\end{table}

Not modelled: \texttt{tokio}, \texttt{async-trait} and \texttt{pingora-core},
the async layer and the process's life cycle; \texttt{io-uring}, the kernel's
side of the ring; and \texttt{pingora-prometheus}, which does not change what
the backend does.

In November 2026, Veil 2.0 pinned Lean 4.35.0, and the Lean library of
Aeneas required Lean 4.36.0 with Mathlib. The file-system library and the
storage contract depend on Lean's core library alone, so they build under
both; this project's own Lean code sits in two projects: the protocol model
under Veil's version, the proofs about the code under that of Aeneas.

\section{When the Checks Run}
\label{app:verification-schedule}

Table~\ref{tab:verification-schedule} gives when each check runs, in jobs
that run side by side. No check has a time limit:
\reqref{HLR-VERIFICATION.FORMAL-01} leaves how long a check may run to a
later decision, and no job sets a limit of its own. Each job records how
long its check took, in the record of checked properties, and that decision
is to be made from those durations. Small bounds run on each pull request,
and deeper ones daily.

\begin{table}[tb]
\centering\small\renewcommand{\arraystretch}{1.3}%
\caption{When the checks run}\label{tab:verification-schedule}
\begin{tabularx}{\linewidth}{|>{\raggedright\arraybackslash}X|>{\raggedright\arraybackslash}p{0.4\linewidth}|}
\hline
\textbf{Check} & \textbf{When} \\
\hline
Kani, small bounds & Each pull request \\
\hline
loom and shuttle, small bounds & Each pull request \\
\hline
The Lean libraries and the protocol model & Each pull request that changes them \\
\hline
Aeneas: the translation up to date, and the proofs about the code & Each pull request that touches a checked module \\
\hline
The crash explorer, small scenarios & Each pull request \\
\hline
Deeper Kani, loom and shuttle bounds; the model's exhaustive checks; the full crash exploration; trace checking & Daily \\
\hline
The \texttt{log-writes} replays, ext4 and XFS & Daily, in a privileged job \\
\hline
Each library against what it models, in its own workflow & Each pull request that changes the library, and daily \\
\hline
\end{tabularx}
\end{table}

Each check runs in a job of its own: the prebuilt files of Mathlib alone
take several gigabytes, and the existing jobs already free disk space before
they build. The translation Aeneas produces is committed, and its check
regenerates it and fails on a difference, as the formatting check does, so
that a reviewer sees which proofs a change reaches. A job whose inputs a
change does not touch finishes at once, but still reports, so that it can be
a required check.

For now no check blocks a merge or a release
(\reqref{ADR-VERIFICATION.FORMAL-01}): the checks report their results, as
the daily runs of the compliance workflow do.

\section{Risks}
\label{app:verification-risks}

\begin{itemize}
  \item \emph{A wrong model.} A proof about a model that does not match what
    it models proves nothing about the code. The comparisons with Linux and
    with the crates, the cross-check with the \texttt{log-writes} replays, and
    trace checking are there to catch it.
  \item \emph{Young tools}: Charon in beta, the function contracts of Kani
    unstable, Veil young, and loom missing some executions by its own
    account.
  \item \emph{Toolchains that move.} The Lean versions, Mathlib, Charon and
    Kani change often; each is pinned, in the Nix shell or the Lake
    manifest, and an upgrade is a reviewed change like any other.
  \item \emph{Upkeep while the code changes.} The modules that have settled
    are checked first, and phase~5 follows the work on the write path.
\end{itemize}

\section{What Stays Unchecked}
\label{app:verification-unchecked}

\begin{itemize}
  \item The async layer and Tokio's scheduling.
  \item Pingora beyond the storage contract.
  \item The kernel and the file systems: their behaviour is assumed, per
    profile, and checked only by the \texttt{log-writes} replays.
  \item A storage device that honours flushes, which
    \reqref{HLR-PERSISTENCE.CRASHSAFE-02} already assumes.
  \item The kernel's side of \texttt{io\_uring}.
\end{itemize}
